\section{Coordinated projections of the HMT state: algebraic equivalence of screens and physical realization in dimension 11}
\label{sec:cierre-equivalencias-moonshine-teoria-m-20260825}

Let \(\mathcal X_{\mathrm{HMT}}^{\mathrm{enr}}\) be the domain preserving
phase, orientation, quotient, carry, boundary, incidence, grading, and
memory. Its exceptional and screen components are extracted by the
canonical projections
\[
 \varepsilon_{\mathrm{exc}}:
 \mathcal X_{\mathrm{HMT}}^{\mathrm{enr}}
 \longrightarrow\mathscr X_{\mathrm{exc}},
 \qquad
 \varepsilon_{\mathrm{scr}}:
 \mathcal X_{\mathrm{HMT}}^{\mathrm{enr}}
 \longrightarrow\mathscr S_{\mathrm{HMT}}.
\]
The former feeds the already proved construction
\(\mathcal P_{\mathrm{exc}}:\mathscr X_{\mathrm{exc}}\to V^\natural\).
The latter will feed the algebraic morphism
\(\mathcal R_M^{\mathrm{alg}}\) constructed below. The joint realization
is
\[
 \mathcal P_{\mathrm{joint}}^M(x)
 =\left(
   \mathcal P_{\mathrm{exc}}\varepsilon_{\mathrm{exc}}(x),
   \mathcal R_M^{\mathrm{alg}}\varepsilon_{\mathrm{scr}}(x)
  \right).
\]
This map preserves the common antecedent and keeps the codomains distinct.

\subsection{Exact equivalence of lattice screens and conditions for the realization of undecadimensional structures}

The real screen starts from
\[
 V_{12}\cong V_1\oplus V_3\oplus V_{3'}\oplus V_5,
 \qquad
 P_{11}=I_{12}-\frac1{12}J_{12}.
\]
The state \(K\) produces \(u_K=P_3P_{11}K\ne0\), the line
\(\ell_K=\mathbb Ru_K\), and
\[
 V_{11}=\ell_K\oplus V_{10},
 \qquad \dim V_{11}=11,
 \qquad \dim V_{10}=10.
\]
The step \(12\to11\) removes the uniform mode, and the step \(11\to10\) removes
the line polarized by \(K\).

The lattice screen belongs to another category. For
\[
 L_{26}=\Lambda_{24}\oplus U\cong II_{25,1}
\]
and a primitive isotropic vector \(e_+\in U\),
\[
 e_+^\perp=\Lambda_{24}\oplus\mathbb Ze_+,
 \qquad
 \boxed{e_+^\perp/\mathbb Ze_+\cong\Lambda_{24}.}
\]
Orthogonality and the quotient eliminate two directions. This theorem is not
a repetition of \(11\to10\) nor an automatic identification with a
physical spacetime.

\subsection{Algebraic screen morphism and quotient equivalence}

Let
\[
 P_{10}:=P_{11}-\frac{u_Ku_K^{\mathsf T}}{\lVert u_K\rVert^2},
 \qquad
 \mathscr G_K:=
 \{(v_{11},v_{10})\in V_{11}\times V_{10}:
             v_{10}=P_{10}v_{11}\}.
\]
The decomposition
\(e_+^\perp=\Lambda_{24}\oplus\mathbb Ze_+\) defines
\[
 q_{24}:e_+^\perp\longrightarrow\Lambda_{24},
 \qquad q_{24}(\lambda+ae_+)=\lambda.
\]
Take
\[
 \mathscr S_{\mathrm{HMT}}
 :=V_{12}\times\mathscr D_0\times e_+^\perp,
 \qquad
 \mathscr S_M^{\mathrm{alg}}
 :=\mathscr G_K\times\mathscr D_0\times\Lambda_{24}.
\]

\begin{definition}[Algebraic screen morphism M]
\label{def:cierre-vtm-rm-algebraico}
For \(d=(m,w,R_0)\in\mathscr D_0\), we define
\[
 \boxed{
 \mathcal R_M^{\mathrm{alg}}(v,d,y)
 :=\bigl((P_{11}v,P_{10}v),d,q_{24}(y)\bigr).}
\]
\end{definition}

\begin{theorem}[Algebraic Equivalence of Quotient Screens]
\label{thm:cierre-vtm-rm-algebraico}
The preceding morphism induces an isomorphism
\[
 \boxed{
 \overline{\mathcal R}_M^{\mathrm{alg}}:
 (V_{12}/V_1)\times\mathscr D_0
 \times(e_+^\perp/\mathbb Ze_+)
 \xrightarrow{\ \cong\ }
 \mathscr S_M^{\mathrm{alg}}.}
\]
If
\[
 \tau_{\mathrm{src}}([v],d,[y])
 =([v],\mathsf T_0d,[y])
\]
and
\[
 \tau_{\mathrm{tar}}((v_{11},v_{10}),d,\lambda)
 =((v_{11},v_{10}),\mathsf T_0d,\lambda),
\]
then
\[
 \overline{\mathcal R}_M^{\mathrm{alg}}\tau_{\mathrm{src}}
 =\tau_{\mathrm{tar}}\overline{\mathcal R}_M^{\mathrm{alg}}.
\]
The compact energy \(E_{R_0}(m,w)\) is preserved under this
intertwining and under reciprocal radii.
\end{theorem}

\begin{proof}
The projector \(P_{11}\) has kernel \(V_1\) and image \(V_{11}\). Since
\(P_{10}P_{11}=P_{10}\), the map
\[
 [v]\longmapsto(P_{11}v,P_{10}v)
\]
is an isomorphism of \(V_{12}/V_1\) onto the graph
\(\mathscr G_K\). The map \(q_{24}\) is surjective and its kernel is
\(\mathbb Ze_+\), hence it induces the isomorphism
\(e_+^\perp/\mathbb Ze_+\cong\Lambda_{24}\). The factor
\(\mathscr D_0\) is transported by the identity. The product of these three
isomorphisms is \(\overline{\mathcal R}_M^{\mathrm{alg}}\).

The projectors and \(q_{24}\) do not act on \(d\), whereas the two
involutions act only through \(\mathsf T_0\). The square commutes
by direct substitution. The energetic invariance is the
\cref{thm:cierre-vtm-tres-formulaciones-t}.
\end{proof}

The mathematical reduction map, and not merely a
correspondence of ranks, is thus constructed. The screen \(11\to10\), the quotient
\(26\to24\), and radius duality appear as distinct factors of the
same morphism.

\subsection{Physical Recognition Undecadimensional}

The interpretation as M-theory is separated by a recognition map.
\[
 \Xi_M:\mathscr S_M^{\mathrm{alg}}
 \longrightarrow\mathcal X_{11}^{\mathrm{phys}},
\]
whose codomain carries at least
\[
 g_{MN},\qquad C_{MNP},\qquad\psi_M,
\]
the action and the supersymmetry transformations. An operator completion includes
\[
 \mathfrak S_M=(\mathcal H_{\mathrm{phys}},H,Q),
 \qquad Q^2=H,
\]
and an isometry
\[
 W:\mathcal H_{\mathrm{HMT}}\longrightarrow\mathcal H_{\mathrm{phys}}
\]
which, in its common domains, satisfies
\[
 W\widehat H=HW,
 \qquad W\widehat Q=QW,
 \qquad \widehat Q^{\,2}=\widehat H.
\]

\begin{theorem}[Physically consistent recognition criterion]
\label{thm:cierre-vtm-reconocimiento-fisico-m}
If \(\Xi_M\) preserves the inner product, transports the direction \(\ell_K\) to the compact direction, intertwines the duality, and satisfies the preceding operator identities, then the compact physical spectrum is invariant under reciprocal radii and the transported supercharge satisfies \(Q^2=H\) on the image of \(W\).
\end{theorem}

\begin{proof}
The interlinking of \(\Xi_M\) with \(\tau_{\mathrm{tar}}\), joined to
\cref{thm:cierre-vtm-rm-algebraico}, transports the unitary involution to the
compact sector. For the supercharge,
\[
 Q^2W=W\widehat Q^{\,2}=W\widehat H=HW.
\]
The preservation of the inner product converts the transported operator into a unitary equivalence on the image.
\end{proof}

The isomorphism
\(\overline{\mathcal R}_M^{\mathrm{alg}}\) is mathematical and has been proved.
Recognition by \(\Xi_M\) is a physical interface conditioned on the
fields, action, gauge constraints, supersymmetry,
oscillators, amplitudes, anomaly control, and low-energy limit of the selected
realization. The physical codomain is not identified with
\(V^\natural\): both receive coordinated information from the same antecedent
through distinct functors.

A literal action of \(\mathbb M\) on digits is not postulated. This sentence is
the typing control placed at the end; it does not replace the positive chain that
reaches the deep action of \(\mathbb M\) in \(V^\natural\).
